\section{Methodology}
A model written in Fortran was used that implemented the linear shallow water equations, with a stress described by \autoref{eq:stress}. The code was initialized at rest and used an Euler step for the first step, then a leapfrog scheme on a Arakawa C-grid as presented by \citet{Doos2022}. The code used a time-step derived from an appropriate Courant number. The code ran for at 80 days, and was saved every 12th hour. The code can be observed in \autoref{ap:code}.

The code used a $H=\SI{3000}{\meter}$, $f_0=\SI{7.27e-5}{\per\second}$ and $\beta=\SI{1.98e-11}{\per\meter\per\second}$. The model used a domain with solid boundaries of the size $L_x=L_y = \SI{7000}{\kilo\meter}$ to represent the Atlantic Ocean (twice as wide for the Pacific Ocean) with $\Delta x=\Delta y = \SI{35}{\kilo\meter}$. Varying the stress was done by using $U_{10} = [\SI{4}{\meter\per\second},\, \SI{5}{\meter\per\second},\, \SI{6}{\meter\per\second}]$, and varying the viscosity was done by $A = [\SI{3e4}{\meter\squared\per\second},\, \SI{1e5}{\meter\squared\per\second},\, \SI{5e5}{\meter\squared\per\second}]$. The harmonic viscosity in the boundary points assumed zero flow outside the domain, leading to energy dissipation. 